% © 2026 Ing. David Jaroš — CC BY-NC-ND 4.0
% Status: linearized frozen-coefficient analysis of the self-consistent
% D-composite scheme; GAP-10T-DYN remains NARROWED (resonant sector open).
\documentclass[11pt,a4paper]{article}
\usepackage{amsmath,amssymb,amsthm}
\usepackage[most]{tcolorbox}
\usepackage[margin=2.5cm]{geometry}
\newtheorem{theorem}{Theorem}
\newtheorem{lemma}{Lemma}
\newtheorem{corollary}{Corollary}
\newcommand{\B}{\mathbb B}
\newcommand{\dd}{\ddagger}
\title{The Self-Consistent $D$-Composite Scheme at Linear Order:\\
Sector Constraint, Symbol Identity, Off-Resonance Flatness,\\
and the Six-Dimensional Resonant Sector}
\author{Ing. David Jaro\v{s}}
\date{26 July 2026}
\begin{document}\maketitle

\begin{abstract}
We begin the audit of the last unanalyzed variational scheme of the minimal
branch: the self-consistent $D$-composite system in which the connection
sits inside the defining jet.  Four exact results are proved at the
linearized, frozen-coefficient level around the affine point.
(i) Imposing $\Theta\in W_L$ defines a consistent Lorentz-real
subsector; necessity for self-consistent solutions is \emph{not} proved
(a constant Hermitian shift at a constant tetrad is a counterexample).
(ii) The frozen symbol $A(s,\lambda)$ of the linearized implicit system
satisfies the exact operator identity $A^3=(\lambda\!\cdot\!s)\,A^2$; its
spectrum is $\{0^{10},\,q^{6}\}$ with $q=\lambda\!\cdot\!s$, with three
nilpotent Jordan blocks on the kernel and $\det(I-A)=(1-q)^6$.
(iii) Off the single resonance hyperplane $q=1$ the unique
$\delta\Theta$-driven solution is the exact gradient; the linearized
tetrad is holonomic and the linearized geometry pullback-flat.
(iv) At $q=1$ there is an exactly six-dimensional space of homogeneous
resonant modes --- $\dim\mathfrak{so}(1,3)$ of them --- and every one is
anholonomic with linearly independent curls.  All linearized anholonomy,
hence all candidate linearized curvature, is confined to this resonant
sector.  The variable-coefficient assembly of resonant modes, the
quadratic action on the resonant sector, on-shell torsion, and the
chain-rule spin current remain open.  All algebraic statements are
verified in \texttt{tools/verify\_dcomposite\_linearized.py}.
\end{abstract}

\section{The self-consistent system}
The scheme couples the covariant jet back to the connection:
\begin{equation}
  E_\mu=\mathcal N_0^{-1/2}\bigl(\partial_\mu\Theta
  +\Omega_\mu[E]\,\Theta+\Theta\,\Omega_\mu[E]^{\dd}\bigr),
  \qquad
  g_{\mu\nu}=\tfrac12\bigl(E_\mu^\sharp E_\nu+E_\nu^\sharp E_\mu\bigr),
  \label{eq:selfcons}
\end{equation}
with $\Omega_\mu[E]$ the spin representation of the Levi--Civita
connection of the tetrad defined by the slice coordinates of $E_\mu$
(torsion-free branch first; anholonomy, not torsion, is the requirement
for curvature).

\begin{lemma}[$W_L$ subsector]\label{lem:sector}
Demanding $E_\mu\in W_L$ for \emph{every} independently chosen Lorentz
connection forces the Hermitian part of $\Theta$ to vanish; hence
$\Theta\in W_L$ defines a consistent Lorentz-real subsector (exact
solve; check L0).  This does \emph{not} prove necessity for
self-consistent solutions, where only $\Omega=\Omega[E]$ occurs:
$\Theta=\Theta_{\rm aff}+H$ with constant Hermitian $H$ has
$\Omega[E]=0$ and $D_\mu\Theta\in W_L$ although $\Theta\notin W_L$
(check L0b).  All results below are statements about the $W_L$
subsector.
\end{lemma}

\section{Linearization and the frozen symbol}
The affine field $\Theta_{\rm aff}=\Theta_0+\sqrt{\mathcal N_0}\,\bar
E_\nu x^\nu$ with $\Theta_0\in W_L$ is an exact solution of
\eqref{eq:selfcons} with $\Omega=0$.  Linearize
$\Theta=\Theta_{\rm aff}+\varepsilon\theta$,
$E=\bar E+\varepsilon F$, freeze coefficients at $x_0$, and pass to the
real-exponential (Laplace-type) symbol $\partial_\nu\to s_\nu$ with
real $s$; for real Fourier modes $\partial_\nu\to ik_\nu$ the invariant
is $q=i\lambda\!\cdot\!k$, so $q=1$ has no real-$k$ solution and the
sector below consists of exponential/evanescent symbol modes.  The system becomes
\begin{equation}
  (I-A(s,\lambda))\,F=\partial\theta',\qquad
  \lambda^a=\text{slice coords of }\Theta_{\rm aff}(x_0)/\sqrt{\mathcal N_0}
  =x_0^a+\theta_0^a,
\end{equation}
where $A$ implements the linearized Levi--Civita connection of the
perturbed tetrad acting on the background through the bimodule.
$\mathcal N_0$ cancels; $A$ is bilinear in $(s,\lambda)$.  Under the
diagonal Lorentz group at the identity-tetrad background the mixed
contraction $q=\lambda^a s_a$ is invariant.

\begin{theorem}[Symbol identity and spectrum]\label{thm:symbol}
$A^3=q\,A^2$ holds as a polynomial identity in all eight variables
(check D2, symbolic).  Generically $\operatorname{rank}A=9$,
$\operatorname{rank}A^2=6$; hence the characteristic polynomial is
$t^{10}(t-q)^6$, the minimal polynomial $t^2(t-q)$, and
$\det(I-A)=(1-q)^6$ (check D3).  The frozen full symbol is modewise
invertible iff $q\neq1$; local or global uniqueness for the
variable-coefficient PDE does not follow from this alone.
\end{theorem}

\begin{theorem}[Off-resonance flatness]\label{thm:offres}
The linearized Levi--Civita connection factors through the curl of the
tetrad perturbation, so $A$ annihilates every exact gradient
(check D1).  Consequently, for $q\neq1$ the unique
$\delta\Theta$-driven solution of $(I-A)F=\partial\theta'$ is
$F=\partial\theta'$ itself: exactly holonomic, with pullback-flat
linearized geometry (check D4).
\end{theorem}

\begin{theorem}[Resonant sector]\label{thm:res}
At $q=1$ the eigenspace $\ker(I-A)$ is exactly six-dimensional
($=\dim\mathfrak{so}(1,3)$).  Every resonant mode is anholonomic and the
six curls are linearly independent (check D5, exact at a generic
resonant point).  All linearized anholonomy of the $D$-composite scheme
is confined to this sector.
\end{theorem}

\begin{corollary}[Where linearized curvature can live]
Since curvature at linear order requires anholonomy, any linearized
gravitational content of the $D$-composite scheme must be assembled from
the exponential symbol modes on the moving surface
$(x_0+\theta_0)\!\cdot\!s=1$; each such mode carries a nonzero
linearized Riemann image (check D5c), but their relation to
real-frequency propagation is open.  Off it, the scheme is flat at linear
order --- but, unlike the exact-gradient restriction, not identically:
the resonant sector is a genuine nontrivial remainder.
\end{corollary}

\section{Answers to the six audit questions at this level}
(1)~Modewise invertibility of the frozen full symbol: yes iff
$q\neq1$; PDE-level uniqueness open.
(2)~Anholonomy: driven solutions are exactly holonomic; anholonomy
exists only in the resonant sector.
(3)~$R\neq0$ without inserting curvature by hand: at linear order only
through resonant modes; whether they assemble into curved solutions of
the variable-coefficient system is open.
(4)~Chain-rule spin current: not yet computed.
(5)~On-shell torsion selection: not yet computed.
(6)~Degeneracy: off resonance the scheme degenerates to the flat sector,
but not identically --- the six-dimensional resonant remainder is the
irreducible core on which the quadratic action must be evaluated next.

\begin{tcolorbox}[colback=orange!6!white,colframe=orange!70!black,
title=Scope and non-claims]
Frozen-coefficient (microlocal) analysis at the affine point.  Not a
solution theory for the variable-coefficient system: $\lambda$ moves
with $x_0$, so the resonance set sweeps spacetime and the constant-
coefficient solvability statements do not globalize automatically.
Nothing here derives the Hilbert--Palatini term, decides on-shell
torsion, or establishes curved solutions.  A discovered-and-fixed index
transposition in the linearized connection (which leaked the symmetric
tetrad part into $\omega^1$) is documented in the verifier as regression
check D1; results computed with the faulty formula (an earlier
$\{0^4,q^9,(q/2)^3\}$ spectrum) are void.
\end{tcolorbox}

\end{document}
